
\documentclass{article}
\usepackage{colortbl}
\usepackage{makecell}
\usepackage{multirow}
\usepackage{supertabular}

\begin{document}

\newcounter{utterance}

\centering \large Interaction Transcript for game `function\_detective', experiment `string\_active\_pair\_checks', episode 164 with Qwen3.6{-}35B{-}A3B{-}FP8{-}without{-}reasoning.
\vspace{24pt}

{ \footnotesize  \setcounter{utterance}{1}
\setlength{\tabcolsep}{0pt}
\begin{supertabular}{c@{$\;$}|p{.15\linewidth}@{}p{.15\linewidth}p{.15\linewidth}p{.15\linewidth}p{.15\linewidth}p{.15\linewidth}}
    \# & \multicolumn{2}{c}{Player} && \multicolumn{2}{c}{Game Master} \\
    \hline

    \theutterance \stepcounter{utterance}  
    & & & \multicolumn{4}{p{0.6\linewidth}}{
        \cellcolor[rgb]{0.9,0.9,0.9}{
            \makecell[{{p{\linewidth}}}]{
                \texttt{\tiny{[P1$\langle$GM]}}
                \texttt{I have a hidden Python function for you to figure out.} \\
\\ 
\texttt{Your job is to infer what it does and then submit a correct Python implementation.} \\
\\ 
\texttt{Function signature:} \\
\texttt{(text: str) {-}> str} \\
\\ 
\texttt{You have 15 turns total.} \\
\\ 
\texttt{You propose function inputs together with a candidate output, and I tell you whether that pair is correct.} \\
\\ 
\texttt{You must make your best SOLVE attempt no later than the second{-}to{-}last turn.} \\
\texttt{Do not wait until the very end.} \\
\\ 
\texttt{You can use these actions:} \\
\texttt{{-} JSON format described at the end of the prompt} \\
\texttt{{-} SOLVE: ```python ... ```} \\
\\ 
\texttt{SOLVE: ```python} \\
\texttt{def solution(<value1>):} \\
\texttt{    ...} \\
\texttt{```} \\
\\ 
\texttt{Do not add import statements in your guess.} \\
\texttt{Anything you might need is already available, so just write the function body and return the answer.} \\
\\ 
\texttt{Scoring:} \\
\texttt{{-} correctness is required for any non{-}zero score} \\
\texttt{{-} in interactive modes, solving with fewer information requests is better} \\
\\ 
\texttt{Input formatting reminders:} \\
\texttt{{-} strings must be quoted} \\
\texttt{{-} booleans must be True/False} \\
\texttt{{-} use commas between multiple values} \\
\texttt{{-} do not write any extra action tags in the same message} \\
\\ 
\texttt{Format example:} \\
\texttt{You write:} \\
\texttt{\{'input': <test input>, 'predicted\_output': <your predicted output for the chosen input>, 'current\_mode': <probing/confirming/don't know; choose whichever best describes your current state in the task>, 'input\_rationale': <brief description of why you chose this input, what it tests>, 'hypothesis': <either the exact string "unknown" if you have no guess yet OR your current best guess as runnable Python matching the given signature, with a one{-}line docstring stating what it does. Never write placeholder code, if you cannot commit to a guess, use "unknown". Match the given signature.} \\
\texttt{e.g.} \\
\texttt{def f(x):} \\
\texttt{	"""Returns the absolute value of x."""} \\
\texttt{	return x if x >= 0 else {-}x} \\
\\ 
\texttt{Use "unknown" if you have no guess yet.>, 'confidence\_score': <score between 1{-}100. 100 means highest confidence in your hypothesis>\}} \\
\\ 
\texttt{I answer:} \\
\texttt{OUTPUT: True} \\
\\ 
\texttt{Output only one valid json per turn.} \\
            }
        }
    }
     \\ \\

    \theutterance \stepcounter{utterance}  
    & \multicolumn{4}{p{0.6\linewidth}}{
        \cellcolor[rgb]{0.561,0.737,0.561}{
            \makecell[{{p{\linewidth}}}]{
                \texttt{\tiny{[P1$\rangle$GM]}}
                \texttt{\{'input': 'hello', 'predicted\_output': 'hello', 'current\_mode': 'probing', 'input\_rationale': 'Starting with a simple lowercase word to see if the function modifies it in any way.', 'hypothesis': 'unknown', 'confidence\_score': 10\}} \\
            }
        }
    }
    & & \\ \\

    \theutterance \stepcounter{utterance}  
    & & & \multicolumn{4}{p{0.6\linewidth}}{
        \cellcolor[rgb]{0.9,0.9,0.9}{
            \makecell[{{p{\linewidth}}}]{
                \texttt{\tiny{[P1$\langle$GM]}}
                \texttt{Game Violation: ParseError: Response is not valid JSON.} \\
\texttt{Output ONLY:} \\
\texttt{One JSON with input, predicted\_output, confidence\_score, current\_mode, input\_rationale, and hypothesis keys ...} \\
\texttt{OR} \\
\texttt{SOLVE: ```python ... ```} \\
            }
        }
    }
     \\ \\

    \theutterance \stepcounter{utterance}  
    & \multicolumn{4}{p{0.6\linewidth}}{
        \cellcolor[rgb]{0.561,0.737,0.561}{
            \makecell[{{p{\linewidth}}}]{
                \texttt{\tiny{[P1$\rangle$GM]}}
                \texttt{\{} \\
\texttt{  "input": "hello",} \\
\texttt{  "predicted\_output": "hello",} \\
\texttt{  "current\_mode": "probing",} \\
\texttt{  "input\_rationale": "Starting with a simple lowercase word to see if the function modifies it in any way.",} \\
\texttt{  "hypothesis": "unknown",} \\
\texttt{  "confidence\_score": 10} \\
\texttt{\}} \\
            }
        }
    }
    & & \\ \\

    \theutterance \stepcounter{utterance}  
    & & & \multicolumn{4}{p{0.6\linewidth}}{
        \cellcolor[rgb]{0.9,0.9,0.9}{
            \makecell[{{p{\linewidth}}}]{
                \texttt{\tiny{[P1$\langle$GM]}}
                \texttt{OUTPUT: False} \\
            }
        }
    }
     \\ \\

    \theutterance \stepcounter{utterance}  
    & \multicolumn{4}{p{0.6\linewidth}}{
        \cellcolor[rgb]{0.561,0.737,0.561}{
            \makecell[{{p{\linewidth}}}]{
                \texttt{\tiny{[P1$\rangle$GM]}}
                \texttt{\{} \\
\texttt{  "input": "hello",} \\
\texttt{  "predicted\_output": "hello",} \\
\texttt{  "current\_mode": "probing",} \\
\texttt{  "input\_rationale": "The previous output was False, meaning 'hello' is not the correct output for input 'hello'. This suggests the function does something to the string. Let's try a single character to see if it's a simple transformation like reversing or changing case.",} \\
\texttt{  "hypothesis": "unknown",} \\
\texttt{  "confidence\_score": 10} \\
\texttt{\}} \\
            }
        }
    }
    & & \\ \\

    \theutterance \stepcounter{utterance}  
    & & & \multicolumn{4}{p{0.6\linewidth}}{
        \cellcolor[rgb]{0.9,0.9,0.9}{
            \makecell[{{p{\linewidth}}}]{
                \texttt{\tiny{[P1$\langle$GM]}}
                \texttt{OUTPUT: False} \\
            }
        }
    }
     \\ \\

    \theutterance \stepcounter{utterance}  
    & \multicolumn{4}{p{0.6\linewidth}}{
        \cellcolor[rgb]{0.561,0.737,0.561}{
            \makecell[{{p{\linewidth}}}]{
                \texttt{\tiny{[P1$\rangle$GM]}}
                \texttt{\{} \\
\texttt{  "input": "a",} \\
\texttt{  "predicted\_output": "a",} \\
\texttt{  "current\_mode": "probing",} \\
\texttt{  "input\_rationale": "Testing a single character to see if the function preserves it or changes it. If 'hello' {-}> False, maybe it's not identity. Let's check if it's case{-}sensitive or something else.",} \\
\texttt{  "hypothesis": "unknown",} \\
\texttt{  "confidence\_score": 10} \\
\texttt{\}} \\
            }
        }
    }
    & & \\ \\

    \theutterance \stepcounter{utterance}  
    & & & \multicolumn{4}{p{0.6\linewidth}}{
        \cellcolor[rgb]{0.9,0.9,0.9}{
            \makecell[{{p{\linewidth}}}]{
                \texttt{\tiny{[P1$\langle$GM]}}
                \texttt{OUTPUT: True} \\
            }
        }
    }
     \\ \\

    \theutterance \stepcounter{utterance}  
    & \multicolumn{4}{p{0.6\linewidth}}{
        \cellcolor[rgb]{0.561,0.737,0.561}{
            \makecell[{{p{\linewidth}}}]{
                \texttt{\tiny{[P1$\rangle$GM]}}
                \texttt{\{} \\
\texttt{  "input": "a",} \\
\texttt{  "predicted\_output": "a",} \\
\texttt{  "current\_mode": "probing",} \\
\texttt{  "input\_rationale": "The function returns 'a' for input 'a'. This suggests it might be an identity function for single characters, or perhaps it reverses strings (which is also identity for length 1). Let's test a two{-}character string to distinguish between identity and reversal.",} \\
\texttt{  "hypothesis": "unknown",} \\
\texttt{  "confidence\_score": 10} \\
\texttt{\}} \\
            }
        }
    }
    & & \\ \\

    \theutterance \stepcounter{utterance}  
    & & & \multicolumn{4}{p{0.6\linewidth}}{
        \cellcolor[rgb]{0.9,0.9,0.9}{
            \makecell[{{p{\linewidth}}}]{
                \texttt{\tiny{[P1$\langle$GM]}}
                \texttt{OUTPUT: True} \\
            }
        }
    }
     \\ \\

    \theutterance \stepcounter{utterance}  
    & \multicolumn{4}{p{0.6\linewidth}}{
        \cellcolor[rgb]{0.561,0.737,0.561}{
            \makecell[{{p{\linewidth}}}]{
                \texttt{\tiny{[P1$\rangle$GM]}}
                \texttt{\{} \\
\texttt{  "input": "ab",} \\
\texttt{  "predicted\_output": "ab",} \\
\texttt{  "current\_mode": "probing",} \\
\texttt{  "input\_rationale": "Testing a two{-}character string. If the function is identity, output is 'ab'. If it's reverse, output is 'ba'. This helps distinguish between simple transformations.",} \\
\texttt{  "hypothesis": "unknown",} \\
\texttt{  "confidence\_score": 10} \\
\texttt{\}} \\
            }
        }
    }
    & & \\ \\

    \theutterance \stepcounter{utterance}  
    & & & \multicolumn{4}{p{0.6\linewidth}}{
        \cellcolor[rgb]{0.9,0.9,0.9}{
            \makecell[{{p{\linewidth}}}]{
                \texttt{\tiny{[P1$\langle$GM]}}
                \texttt{OUTPUT: True} \\
            }
        }
    }
     \\ \\

    \theutterance \stepcounter{utterance}  
    & \multicolumn{4}{p{0.6\linewidth}}{
        \cellcolor[rgb]{0.561,0.737,0.561}{
            \makecell[{{p{\linewidth}}}]{
                \texttt{\tiny{[P1$\rangle$GM]}}
                \texttt{\{} \\
\texttt{  "input": "ab",} \\
\texttt{  "predicted\_output": "ab",} \\
\texttt{  "current\_mode": "probing",} \\
\texttt{  "input\_rationale": "The function returns 'ab' for input 'ab'. This is consistent with an identity function. Let's test a longer string to confirm.",} \\
\texttt{  "hypothesis": "unknown",} \\
\texttt{  "confidence\_score": 10} \\
\texttt{\}} \\
            }
        }
    }
    & & \\ \\

    \theutterance \stepcounter{utterance}  
    & & & \multicolumn{4}{p{0.6\linewidth}}{
        \cellcolor[rgb]{0.9,0.9,0.9}{
            \makecell[{{p{\linewidth}}}]{
                \texttt{\tiny{[P1$\langle$GM]}}
                \texttt{OUTPUT: True} \\
            }
        }
    }
     \\ \\

    \theutterance \stepcounter{utterance}  
    & \multicolumn{4}{p{0.6\linewidth}}{
        \cellcolor[rgb]{0.561,0.737,0.561}{
            \makecell[{{p{\linewidth}}}]{
                \texttt{\tiny{[P1$\rangle$GM]}}
                \texttt{\{} \\
\texttt{  "input": "hello world",} \\
\texttt{  "predicted\_output": "hello world",} \\
\texttt{  "current\_mode": "probing",} \\
\texttt{  "input\_rationale": "Testing a string with a space to ensure spaces are preserved and the function is truly identity.",} \\
\texttt{  "hypothesis": "unknown",} \\
\texttt{  "confidence\_score": 10} \\
\texttt{\}} \\
            }
        }
    }
    & & \\ \\

    \theutterance \stepcounter{utterance}  
    & & & \multicolumn{4}{p{0.6\linewidth}}{
        \cellcolor[rgb]{0.9,0.9,0.9}{
            \makecell[{{p{\linewidth}}}]{
                \texttt{\tiny{[P1$\langle$GM]}}
                \texttt{OUTPUT: False} \\
            }
        }
    }
     \\ \\

    \theutterance \stepcounter{utterance}  
    & \multicolumn{4}{p{0.6\linewidth}}{
        \cellcolor[rgb]{0.561,0.737,0.561}{
            \makecell[{{p{\linewidth}}}]{
                \texttt{\tiny{[P1$\rangle$GM]}}
                \texttt{\{} \\
\texttt{  "input": "hello world",} \\
\texttt{  "predicted\_output": "hello world",} \\
\texttt{  "current\_mode": "probing",} \\
\texttt{  "input\_rationale": "The function returned False for 'hello world', meaning the output is not 'hello world'. Since it was identity for 'a' and 'ab', but not for 'hello world', it likely involves some transformation that doesn't affect short strings or specific patterns. Let's try a string with uppercase letters to check for case conversion.",} \\
\texttt{  "hypothesis": "unknown",} \\
\texttt{  "confidence\_score": 10} \\
\texttt{\}} \\
            }
        }
    }
    & & \\ \\

    \theutterance \stepcounter{utterance}  
    & & & \multicolumn{4}{p{0.6\linewidth}}{
        \cellcolor[rgb]{0.9,0.9,0.9}{
            \makecell[{{p{\linewidth}}}]{
                \texttt{\tiny{[P1$\langle$GM]}}
                \texttt{OUTPUT: False} \\
            }
        }
    }
     \\ \\

    \theutterance \stepcounter{utterance}  
    & \multicolumn{4}{p{0.6\linewidth}}{
        \cellcolor[rgb]{0.561,0.737,0.561}{
            \makecell[{{p{\linewidth}}}]{
                \texttt{\tiny{[P1$\rangle$GM]}}
                \texttt{\{} \\
\texttt{  "input": "Hello",} \\
\texttt{  "predicted\_output": "Hello",} \\
\texttt{  "current\_mode": "probing",} \\
\texttt{  "input\_rationale": "Testing a string with an uppercase letter to see if case is preserved or changed. Previous tests with lowercase 'a', 'ab', 'hello' (wait, 'hello' was False? No, 'hello' was False for output 'hello', meaning output != 'hello'. 'a' and 'ab' were True for output 'a' and 'ab' respectively. 'hello world' was False for output 'hello world'. This is confusing. Let's re{-}evaluate. 'hello' {-}> False (output != 'hello'). 'a' {-}> True (output == 'a'). 'ab' {-}> True (output == 'ab'). 'hello world' {-}> False (output != 'hello world'). This suggests the function might be identity for strings of length 1 or 2, but not for longer ones? Or maybe it's about specific characters. Let's try 'abc' to see if length 3 breaks it.",} \\
\texttt{  "hypothesis": "unknown",} \\
\texttt{  "confidence\_score": 10} \\
\texttt{\}} \\
            }
        }
    }
    & & \\ \\

    \theutterance \stepcounter{utterance}  
    & & & \multicolumn{4}{p{0.6\linewidth}}{
        \cellcolor[rgb]{0.9,0.9,0.9}{
            \makecell[{{p{\linewidth}}}]{
                \texttt{\tiny{[P1$\langle$GM]}}
                \texttt{OUTPUT: False} \\
            }
        }
    }
     \\ \\

    \theutterance \stepcounter{utterance}  
    & \multicolumn{4}{p{0.6\linewidth}}{
        \cellcolor[rgb]{0.561,0.737,0.561}{
            \makecell[{{p{\linewidth}}}]{
                \texttt{\tiny{[P1$\rangle$GM]}}
                \texttt{\{} \\
\texttt{  "input": "abc",} \\
\texttt{  "predicted\_output": "abc",} \\
\texttt{  "current\_mode": "probing",} \\
\texttt{  "input\_rationale": "Testing a 3{-}character string. 'a' and 'ab' were identity, but 'hello' (5 chars) and 'hello world' (11 chars) were not. 'abc' is 3 chars. If this is also not identity, it suggests the function is only identity for length <= 2. If it is identity, then 'hello' being non{-}identity is strange unless it's not just length{-}based. Let's check if 'abc' returns 'abc'.",} \\
\texttt{  "hypothesis": "unknown",} \\
\texttt{  "confidence\_score": 10} \\
\texttt{\}} \\
            }
        }
    }
    & & \\ \\

    \theutterance \stepcounter{utterance}  
    & & & \multicolumn{4}{p{0.6\linewidth}}{
        \cellcolor[rgb]{0.9,0.9,0.9}{
            \makecell[{{p{\linewidth}}}]{
                \texttt{\tiny{[P1$\langle$GM]}}
                \texttt{OUTPUT: True} \\
            }
        }
    }
     \\ \\

    \theutterance \stepcounter{utterance}  
    & \multicolumn{4}{p{0.6\linewidth}}{
        \cellcolor[rgb]{0.561,0.737,0.561}{
            \makecell[{{p{\linewidth}}}]{
                \texttt{\tiny{[P1$\rangle$GM]}}
                \texttt{\{} \\
\texttt{  "input": "abcd",} \\
\texttt{  "predicted\_output": "abcd",} \\
\texttt{  "current\_mode": "probing",} \\
\texttt{  "input\_rationale": "Testing a 4{-}character string. 'a', 'ab', 'abc' were identity. 'hello' (5 chars) was not. This suggests the function might be identity for strings of length <= 3, but not for length >= 4? Or perhaps it's related to specific characters. Let's test 'abcd' to see if length 4 breaks the pattern.",} \\
\texttt{  "hypothesis": "unknown",} \\
\texttt{  "confidence\_score": 10} \\
\texttt{\}} \\
            }
        }
    }
    & & \\ \\

    \theutterance \stepcounter{utterance}  
    & & & \multicolumn{4}{p{0.6\linewidth}}{
        \cellcolor[rgb]{0.9,0.9,0.9}{
            \makecell[{{p{\linewidth}}}]{
                \texttt{\tiny{[P1$\langle$GM]}}
                \texttt{OUTPUT: False} \\
            }
        }
    }
     \\ \\

    \theutterance \stepcounter{utterance}  
    & \multicolumn{4}{p{0.6\linewidth}}{
        \cellcolor[rgb]{0.561,0.737,0.561}{
            \makecell[{{p{\linewidth}}}]{
                \texttt{\tiny{[P1$\rangle$GM]}}
                \texttt{\{} \\
\texttt{  "input": "abcd",} \\
\texttt{  "predicted\_output": "abcd",} \\
\texttt{  "current\_mode": "probing",} \\
\texttt{  "input\_rationale": "Testing a 4{-}character string. 'a' (len 1), 'ab' (len 2), 'abc' (len 3) were identity. 'hello' (len 5) and 'hello world' (len 11) were not. 'abcd' (len 4) is the next step. If it's not identity, the function is identity for length <= 3. If it is identity, then 'hello' being non{-}identity is puzzling. Let's see.",} \\
\texttt{  "hypothesis": "unknown",} \\
\texttt{  "confidence\_score": 10} \\
\texttt{\}} \\
            }
        }
    }
    & & \\ \\

    \theutterance \stepcounter{utterance}  
    & & & \multicolumn{4}{p{0.6\linewidth}}{
        \cellcolor[rgb]{0.9,0.9,0.9}{
            \makecell[{{p{\linewidth}}}]{
                \texttt{\tiny{[P1$\langle$GM]}}
                \texttt{OUTPUT: False} \\
            }
        }
    }
     \\ \\

    \theutterance \stepcounter{utterance}  
    & \multicolumn{4}{p{0.6\linewidth}}{
        \cellcolor[rgb]{0.561,0.737,0.561}{
            \makecell[{{p{\linewidth}}}]{
                \texttt{\tiny{[P1$\rangle$GM]}}
                \texttt{\{} \\
\texttt{  "input": "abcde",} \\
\texttt{  "predicted\_output": "abcde",} \\
\texttt{  "current\_mode": "probing",} \\
\texttt{  "input\_rationale": "Testing a 5{-}character string 'abcde'. 'hello' (5 chars) was not identity. If 'abcde' is also not identity, it confirms length >= 4 or 5 breaks identity. If it is identity, then 'hello' has some special property. Let's check if 'abcde' is identity.",} \\
\texttt{  "hypothesis": "unknown",} \\
\texttt{  "confidence\_score": 10} \\
\texttt{\}} \\
            }
        }
    }
    & & \\ \\

    \theutterance \stepcounter{utterance}  
    & & & \multicolumn{4}{p{0.6\linewidth}}{
        \cellcolor[rgb]{0.9,0.9,0.9}{
            \makecell[{{p{\linewidth}}}]{
                \texttt{\tiny{[P1$\langle$GM]}}
                \texttt{OUTPUT: False} \\
            }
        }
    }
     \\ \\

    \theutterance \stepcounter{utterance}  
    & \multicolumn{4}{p{0.6\linewidth}}{
        \cellcolor[rgb]{0.561,0.737,0.561}{
            \makecell[{{p{\linewidth}}}]{
                \texttt{\tiny{[P1$\rangle$GM]}}
                \texttt{\{} \\
\texttt{  "input": "abcde",} \\
\texttt{  "predicted\_output": "abcde",} \\
\texttt{  "current\_mode": "probing",} \\
\texttt{  "input\_rationale": "Testing a 5{-}character string 'abcde'. 'hello' (5 chars) was not identity. If 'abcde' is also not identity, it confirms that strings of length >= 4 are not identity. This strongly suggests the function is identity for strings of length <= 3, and does something else for longer strings. Let's try to figure out what it does for longer strings. Maybe it reverses them? 'hello' reversed is 'olleh'. Let's test if the output for 'hello' is 'olleh'.",} \\
\texttt{  "hypothesis": "unknown",} \\
\texttt{  "confidence\_score": 10} \\
\texttt{\}} \\
            }
        }
    }
    & & \\ \\

    \theutterance \stepcounter{utterance}  
    & & & \multicolumn{2}{p{0.3\linewidth}}{
        \cellcolor[rgb]{0.9,0.9,0.9}{
            \makecell[{{p{\linewidth}}}]{
                \texttt{\tiny{[GM$|$GM]}}
                \texttt{Game Verdict: LOSS (turn limit reached)} \\
            }
        }
    }
    & & \\ \\

\end{supertabular}
}

\end{document}
